\documentclass[11pt,a4paper]{article}
% ============================================
% GIFT v3.4 — Shared Preamble
% ============================================
% Usage: % ============================================
% GIFT v3.4 — Shared Preamble
% ============================================
% Usage: % ============================================
% GIFT v3.4 — Shared Preamble
% ============================================
% Usage: \input{gift_preamble} after \documentclass

% ENCODING & FONTS
\usepackage[utf8]{inputenc}
\usepackage[T1]{fontenc}
\usepackage{lmodern}

% PAGE LAYOUT (Golden Ratio)
\usepackage[margin=1.618cm, top=2.618cm, bottom=2.618cm]{geometry}

% ESSENTIAL PACKAGES
\usepackage{float}
\usepackage{caption}
\usepackage{subcaption}
\usepackage{setspace}
\usepackage{fancyhdr}
\usepackage{xcolor}
\usepackage{hyperref}
\usepackage{csquotes}
\usepackage{amsmath}
\usepackage{amssymb}
\usepackage{amsthm}
\usepackage{mathtools}
\usepackage{booktabs}
\usepackage{longtable}
\usepackage{array}
\usepackage{multirow}
\usepackage{tikz}
\usepackage{graphicx}
\usepackage{listings}
\usepackage{enumitem}
% Unicode fallbacks (pandoc artifacts that survive markdown→LaTeX)
\DeclareUnicodeCharacter{00B0}{\ensuremath{^\circ}}
\DeclareUnicodeCharacter{221A}{\ensuremath{\sqrt{\,}}}
\DeclareUnicodeCharacter{2713}{\ensuremath{\checkmark}}
\DeclareUnicodeCharacter{2717}{\ensuremath{\times}}
\DeclareUnicodeCharacter{2070}{\ensuremath{^{0}}}
\DeclareUnicodeCharacter{00B9}{\ensuremath{^{1}}}
\DeclareUnicodeCharacter{00B2}{\ensuremath{^{2}}}
\DeclareUnicodeCharacter{00B3}{\ensuremath{^{3}}}
\DeclareUnicodeCharacter{2074}{\ensuremath{^{4}}}
\DeclareUnicodeCharacter{2075}{\ensuremath{^{5}}}
\DeclareUnicodeCharacter{2076}{\ensuremath{^{6}}}
\DeclareUnicodeCharacter{2077}{\ensuremath{^{7}}}
\DeclareUnicodeCharacter{2078}{\ensuremath{^{8}}}
\DeclareUnicodeCharacter{2079}{\ensuremath{^{9}}}
\DeclareUnicodeCharacter{2071}{\ensuremath{^{i}}}
\DeclareUnicodeCharacter{2080}{\ensuremath{_{0}}}
\DeclareUnicodeCharacter{2081}{\ensuremath{_{1}}}
\DeclareUnicodeCharacter{2082}{\ensuremath{_{2}}}
\DeclareUnicodeCharacter{2083}{\ensuremath{_{3}}}
\DeclareUnicodeCharacter{2084}{\ensuremath{_{4}}}
\DeclareUnicodeCharacter{2085}{\ensuremath{_{5}}}
\DeclareUnicodeCharacter{2086}{\ensuremath{_{6}}}
\DeclareUnicodeCharacter{2087}{\ensuremath{_{7}}}
\DeclareUnicodeCharacter{2088}{\ensuremath{_{8}}}
\DeclareUnicodeCharacter{2089}{\ensuremath{_{9}}}
\DeclareUnicodeCharacter{208A}{\ensuremath{_{+}}}
\DeclareUnicodeCharacter{208B}{\ensuremath{_{-}}}
\DeclareUnicodeCharacter{2192}{\ensuremath{\to}}
\DeclareUnicodeCharacter{2502}{|}
\DeclareUnicodeCharacter{25BC}{\ensuremath{\blacktriangledown}}
\DeclareUnicodeCharacter{2500}{\ensuremath{-}}
\DeclareUnicodeCharacter{2514}{\ensuremath{\llcorner}}
\DeclareUnicodeCharacter{251C}{\ensuremath{\vdash}}
\DeclareUnicodeCharacter{2026}{\ldots}
\DeclareUnicodeCharacter{2208}{\ensuremath{\in}}
\DeclareUnicodeCharacter{2229}{\ensuremath{\cap}}
\DeclareUnicodeCharacter{222A}{\ensuremath{\cup}}
\DeclareUnicodeCharacter{2295}{\ensuremath{\oplus}}
\DeclareUnicodeCharacter{2297}{\ensuremath{\otimes}}
\DeclareUnicodeCharacter{2245}{\ensuremath{\cong}}
\DeclareUnicodeCharacter{2248}{\ensuremath{\approx}}
\DeclareUnicodeCharacter{2260}{\ensuremath{\neq}}
\DeclareUnicodeCharacter{2261}{\ensuremath{\equiv}}
\DeclareUnicodeCharacter{2264}{\ensuremath{\leq}}
\DeclareUnicodeCharacter{2265}{\ensuremath{\geq}}
\DeclareUnicodeCharacter{00D7}{\ensuremath{\times}}
\DeclareUnicodeCharacter{00B1}{\ensuremath{\pm}}
\DeclareUnicodeCharacter{2225}{\ensuremath{\|}}
% Greek letters used in body text (pandoc passes them through unchanged)
\DeclareUnicodeCharacter{03B1}{\ensuremath{\alpha}}
\DeclareUnicodeCharacter{03B2}{\ensuremath{\beta}}
\DeclareUnicodeCharacter{03B3}{\ensuremath{\gamma}}
\DeclareUnicodeCharacter{03B4}{\ensuremath{\delta}}
\DeclareUnicodeCharacter{03B5}{\ensuremath{\varepsilon}}
\DeclareUnicodeCharacter{03B6}{\ensuremath{\zeta}}
\DeclareUnicodeCharacter{03B7}{\ensuremath{\eta}}
\DeclareUnicodeCharacter{03B8}{\ensuremath{\theta}}
\DeclareUnicodeCharacter{03B9}{\ensuremath{\iota}}
\DeclareUnicodeCharacter{03BA}{\ensuremath{\kappa}}
\DeclareUnicodeCharacter{03BB}{\ensuremath{\lambda}}
\DeclareUnicodeCharacter{03BC}{\ensuremath{\mu}}
\DeclareUnicodeCharacter{03BD}{\ensuremath{\nu}}
\DeclareUnicodeCharacter{03BE}{\ensuremath{\xi}}
\DeclareUnicodeCharacter{03C0}{\ensuremath{\pi}}
\DeclareUnicodeCharacter{03C1}{\ensuremath{\rho}}
\DeclareUnicodeCharacter{03C3}{\ensuremath{\sigma}}
\DeclareUnicodeCharacter{03C4}{\ensuremath{\tau}}
\DeclareUnicodeCharacter{03C5}{\ensuremath{\upsilon}}
\DeclareUnicodeCharacter{03C6}{\ensuremath{\varphi}}
\DeclareUnicodeCharacter{03C7}{\ensuremath{\chi}}
\DeclareUnicodeCharacter{03C8}{\ensuremath{\psi}}
\DeclareUnicodeCharacter{03C9}{\ensuremath{\omega}}
\DeclareUnicodeCharacter{0394}{\ensuremath{\Delta}}
\DeclareUnicodeCharacter{0398}{\ensuremath{\Theta}}
\DeclareUnicodeCharacter{039B}{\ensuremath{\Lambda}}
\DeclareUnicodeCharacter{03A0}{\ensuremath{\Pi}}
\DeclareUnicodeCharacter{03A3}{\ensuremath{\Sigma}}
\DeclareUnicodeCharacter{03A6}{\ensuremath{\Phi}}
\DeclareUnicodeCharacter{03A8}{\ensuremath{\Psi}}
\DeclareUnicodeCharacter{03A9}{\ensuremath{\Omega}}
\DeclareUnicodeCharacter{2200}{\ensuremath{\forall}}
\DeclareUnicodeCharacter{2203}{\ensuremath{\exists}}
\DeclareUnicodeCharacter{2207}{\ensuremath{\nabla}}
\DeclareUnicodeCharacter{2202}{\ensuremath{\partial}}
\DeclareUnicodeCharacter{2218}{\ensuremath{\circ}}
\DeclareUnicodeCharacter{2032}{\ensuremath{\prime}}
\DeclareUnicodeCharacter{2113}{\ensuremath{\ell}}
\DeclareUnicodeCharacter{210F}{\ensuremath{\hbar}}
\DeclareUnicodeCharacter{27E8}{\ensuremath{\langle}}
\DeclareUnicodeCharacter{27E9}{\ensuremath{\rangle}}
\DeclareUnicodeCharacter{2308}{\ensuremath{\lceil}}
\DeclareUnicodeCharacter{2309}{\ensuremath{\rceil}}
\DeclareUnicodeCharacter{230A}{\ensuremath{\lfloor}}
\DeclareUnicodeCharacter{230B}{\ensuremath{\rfloor}}
\DeclareUnicodeCharacter{222B}{\ensuremath{\int}}
\DeclareUnicodeCharacter{2211}{\ensuremath{\sum}}
\DeclareUnicodeCharacter{220F}{\ensuremath{\prod}}
\DeclareUnicodeCharacter{221E}{\ensuremath{\infty}}
\DeclareUnicodeCharacter{00B5}{\ensuremath{\mu}}
\DeclareUnicodeCharacter{2236}{\ensuremath{\colon}}
\DeclareUnicodeCharacter{2192}{\ensuremath{\to}}
\DeclareUnicodeCharacter{21A6}{\ensuremath{\mapsto}}
\DeclareUnicodeCharacter{21D2}{\ensuremath{\Rightarrow}}
\DeclareUnicodeCharacter{21D4}{\ensuremath{\Leftrightarrow}}
\DeclareUnicodeCharacter{2282}{\ensuremath{\subset}}
\DeclareUnicodeCharacter{2286}{\ensuremath{\subseteq}}
\DeclareUnicodeCharacter{2287}{\ensuremath{\supseteq}}
\DeclareUnicodeCharacter{222B}{\ensuremath{\int}}
\DeclareUnicodeCharacter{2205}{\ensuremath{\emptyset}}
\DeclareUnicodeCharacter{2299}{\ensuremath{\odot}}
\DeclareUnicodeCharacter{2220}{\ensuremath{\angle}}
\DeclareUnicodeCharacter{2135}{\ensuremath{\aleph}}
\DeclareUnicodeCharacter{207B}{\ensuremath{^{-}}}
\DeclareUnicodeCharacter{207A}{\ensuremath{^{+}}}
\DeclareUnicodeCharacter{220E}{\ensuremath{\blacksquare}}
\DeclareUnicodeCharacter{2061}{}
\DeclareUnicodeCharacter{2062}{}
\DeclareUnicodeCharacter{2063}{}
\DeclareUnicodeCharacter{2009}{\,}
\DeclareUnicodeCharacter{200A}{\,}
\DeclareUnicodeCharacter{200B}{}
\DeclareUnicodeCharacter{2212}{\ensuremath{-}}
\DeclareUnicodeCharacter{2194}{\ensuremath{\leftrightarrow}}
\DeclareUnicodeCharacter{2099}{\ensuremath{_{n}}}
\DeclareUnicodeCharacter{2098}{\ensuremath{_{m}}}
\DeclareUnicodeCharacter{2096}{\ensuremath{_{k}}}
\DeclareUnicodeCharacter{2095}{\ensuremath{_{h}}}
\DeclareUnicodeCharacter{2090}{\ensuremath{_{a}}}
\DeclareUnicodeCharacter{2091}{\ensuremath{_{e}}}
\DeclareUnicodeCharacter{2093}{\ensuremath{_{x}}}
\DeclareUnicodeCharacter{1D62}{\ensuremath{_{i}}}
\DeclareUnicodeCharacter{2C7C}{\ensuremath{_{j}}}
\DeclareUnicodeCharacter{2C7B}{\ensuremath{^{e}}}
\DeclareUnicodeCharacter{1D52}{\ensuremath{^{o}}}
\DeclareUnicodeCharacter{1D5}{\ensuremath{^{u}}}
\DeclareUnicodeCharacter{02E1}{\ensuremath{^{l}}}
\DeclareUnicodeCharacter{2295}{\ensuremath{\oplus}}
\DeclareUnicodeCharacter{27FA}{\ensuremath{\Longleftrightarrow}}
\DeclareUnicodeCharacter{27F9}{\ensuremath{\Longrightarrow}}
\DeclareUnicodeCharacter{2299}{\ensuremath{\odot}}
% Provide \tightlist used by pandoc-generated lists (no-op)
\providecommand{\tightlist}{\setlength{\itemsep}{0pt}\setlength{\parskip}{0pt}}
% pandoc longtable column specs: provide \real as passthrough and a
% `none` counter so that calc's \real{x} expansion path works.
\newcounter{none}
\providecommand{\real}[1]{#1}

% LISTINGS
\lstset{
    basicstyle=\small\ttfamily,
    breaklines=true, frame=single,
    keepspaces=true, showstringspaces=false,
    aboveskip=0.8em, belowskip=0.8em
}

% HEADER/FOOTER
\setlength{\headheight}{14pt}
\pagestyle{fancy}
\fancyhf{}
\fancyhead[R]{\thepage}
\renewcommand{\headrulewidth}{0.2pt}

% HYPERREF
\hypersetup{
    colorlinks=true,
    linkcolor=blue, citecolor=blue, urlcolor=blue,
    pdfauthor={Brieuc de La Fourniere}
}

% SPACING
\setstretch{1.2}
\setlength{\parskip}{0.4em}
\setlength{\parindent}{0pt}

% TITLE FORMATTING
\usepackage{titling}
\pretitle{\LARGE\bfseries}
\posttitle{\vspace{-0.4em}}
\preauthor{}
\postauthor{}
\predate{}
\postdate{}
\setlength{\droptitle}{-2.0em}

% CUSTOM COMMANDS — GIFT notation
\newcommand{\E}{\mathrm{E}}
\newcommand{\Gtwo}{\mathrm{G}_2}
\newcommand{\Kseven}{K_7}
\newcommand{\AdS}{\mathrm{AdS}}
\newcommand{\SU}{\mathrm{SU}}
\newcommand{\SO}{\mathrm{SO}}
\newcommand{\U}{\mathrm{U}}
\newcommand{\Sp}{\mathrm{Sp}}
\newcommand{\PSL}{\mathrm{PSL}}
\newcommand{\SM}{\mathrm{SM}}
\newcommand{\Tr}{\mathrm{Tr}}
\newcommand{\rank}{\mathrm{rank}}
\newcommand{\Vol}{\mathrm{Vol}}
\newcommand{\Hol}{\mathrm{Hol}}
\newcommand{\Ric}{\mathrm{Ric}}
\newcommand{\Rm}{\mathrm{Rm}}
\newcommand{\Scal}{\mathrm{Scal}}
\DeclareMathOperator{\diag}{diag}
\DeclareMathOperator{\cond}{cond}
\newcommand{\GIFT}{\textrm{GIFT}}
\newcommand{\CP}{\mathrm{CP}}
\newcommand{\CKM}{\mathrm{CKM}}
\newcommand{\PMNS}{\mathrm{PMNS}}
\newcommand{\EW}{\mathrm{EW}}
\newcommand{\Pl}{\mathrm{Pl}}
\newcommand{\DE}{\mathrm{DE}}
\newcommand{\DM}{\mathrm{DM}}
\newcommand{\KK}{\mathrm{KK}}
\newcommand{\NK}{\mathrm{NK}}
\newcommand{\TCS}{\mathrm{TCS}}
\newcommand{\PINN}{\mathrm{PINN}}
\newcommand{\btwo}{b_2}
\newcommand{\bthree}{b_3}
\newcommand{\typeI}{\textsc{I}}
\newcommand{\typeII}{\textsc{II}}
\newcommand{\typeIII}{\textsc{III}}
\newcommand{\typeIV}{\textsc{IV}}
\newcommand{\proven}{\textsc{Proven}}

% PDF bookmark safety
\pdfstringdefDisableCommands{%
  \def\textsubscript#1{#1}%
  \def\textsuperscript#1{#1}%
  \def\CP{CP}%
  \def\Gtwo{G2}%
  \def\Kseven{K7}%
  \def\GIFT{GIFT}%
  \def\E{E}%
  \def\SM{SM}%
}

% THEOREM ENVIRONMENTS
\newtheorem{theorem}{Theorem}[section]
\newtheorem{proposition}[theorem]{Proposition}
\newtheorem{lemma}[theorem]{Lemma}
\newtheorem{corollary}[theorem]{Corollary}
\theoremstyle{definition}
\newtheorem{definition}[theorem]{Definition}
\theoremstyle{remark}
\newtheorem{remark}[theorem]{Remark}
 after \documentclass

% ENCODING & FONTS
\usepackage[utf8]{inputenc}
\usepackage[T1]{fontenc}
\usepackage{lmodern}

% PAGE LAYOUT (Golden Ratio)
\usepackage[margin=1.618cm, top=2.618cm, bottom=2.618cm]{geometry}

% ESSENTIAL PACKAGES
\usepackage{float}
\usepackage{caption}
\usepackage{subcaption}
\usepackage{setspace}
\usepackage{fancyhdr}
\usepackage{xcolor}
\usepackage{hyperref}
\usepackage{csquotes}
\usepackage{amsmath}
\usepackage{amssymb}
\usepackage{amsthm}
\usepackage{mathtools}
\usepackage{booktabs}
\usepackage{longtable}
\usepackage{array}
\usepackage{multirow}
\usepackage{tikz}
\usepackage{graphicx}
\usepackage{listings}
\usepackage{enumitem}
% Unicode fallbacks (pandoc artifacts that survive markdown→LaTeX)
\DeclareUnicodeCharacter{00B0}{\ensuremath{^\circ}}
\DeclareUnicodeCharacter{221A}{\ensuremath{\sqrt{\,}}}
\DeclareUnicodeCharacter{2713}{\ensuremath{\checkmark}}
\DeclareUnicodeCharacter{2717}{\ensuremath{\times}}
\DeclareUnicodeCharacter{2070}{\ensuremath{^{0}}}
\DeclareUnicodeCharacter{00B9}{\ensuremath{^{1}}}
\DeclareUnicodeCharacter{00B2}{\ensuremath{^{2}}}
\DeclareUnicodeCharacter{00B3}{\ensuremath{^{3}}}
\DeclareUnicodeCharacter{2074}{\ensuremath{^{4}}}
\DeclareUnicodeCharacter{2075}{\ensuremath{^{5}}}
\DeclareUnicodeCharacter{2076}{\ensuremath{^{6}}}
\DeclareUnicodeCharacter{2077}{\ensuremath{^{7}}}
\DeclareUnicodeCharacter{2078}{\ensuremath{^{8}}}
\DeclareUnicodeCharacter{2079}{\ensuremath{^{9}}}
\DeclareUnicodeCharacter{2071}{\ensuremath{^{i}}}
\DeclareUnicodeCharacter{2080}{\ensuremath{_{0}}}
\DeclareUnicodeCharacter{2081}{\ensuremath{_{1}}}
\DeclareUnicodeCharacter{2082}{\ensuremath{_{2}}}
\DeclareUnicodeCharacter{2083}{\ensuremath{_{3}}}
\DeclareUnicodeCharacter{2084}{\ensuremath{_{4}}}
\DeclareUnicodeCharacter{2085}{\ensuremath{_{5}}}
\DeclareUnicodeCharacter{2086}{\ensuremath{_{6}}}
\DeclareUnicodeCharacter{2087}{\ensuremath{_{7}}}
\DeclareUnicodeCharacter{2088}{\ensuremath{_{8}}}
\DeclareUnicodeCharacter{2089}{\ensuremath{_{9}}}
\DeclareUnicodeCharacter{208A}{\ensuremath{_{+}}}
\DeclareUnicodeCharacter{208B}{\ensuremath{_{-}}}
\DeclareUnicodeCharacter{2192}{\ensuremath{\to}}
\DeclareUnicodeCharacter{2502}{|}
\DeclareUnicodeCharacter{25BC}{\ensuremath{\blacktriangledown}}
\DeclareUnicodeCharacter{2500}{\ensuremath{-}}
\DeclareUnicodeCharacter{2514}{\ensuremath{\llcorner}}
\DeclareUnicodeCharacter{251C}{\ensuremath{\vdash}}
\DeclareUnicodeCharacter{2026}{\ldots}
\DeclareUnicodeCharacter{2208}{\ensuremath{\in}}
\DeclareUnicodeCharacter{2229}{\ensuremath{\cap}}
\DeclareUnicodeCharacter{222A}{\ensuremath{\cup}}
\DeclareUnicodeCharacter{2295}{\ensuremath{\oplus}}
\DeclareUnicodeCharacter{2297}{\ensuremath{\otimes}}
\DeclareUnicodeCharacter{2245}{\ensuremath{\cong}}
\DeclareUnicodeCharacter{2248}{\ensuremath{\approx}}
\DeclareUnicodeCharacter{2260}{\ensuremath{\neq}}
\DeclareUnicodeCharacter{2261}{\ensuremath{\equiv}}
\DeclareUnicodeCharacter{2264}{\ensuremath{\leq}}
\DeclareUnicodeCharacter{2265}{\ensuremath{\geq}}
\DeclareUnicodeCharacter{00D7}{\ensuremath{\times}}
\DeclareUnicodeCharacter{00B1}{\ensuremath{\pm}}
\DeclareUnicodeCharacter{2225}{\ensuremath{\|}}
% Greek letters used in body text (pandoc passes them through unchanged)
\DeclareUnicodeCharacter{03B1}{\ensuremath{\alpha}}
\DeclareUnicodeCharacter{03B2}{\ensuremath{\beta}}
\DeclareUnicodeCharacter{03B3}{\ensuremath{\gamma}}
\DeclareUnicodeCharacter{03B4}{\ensuremath{\delta}}
\DeclareUnicodeCharacter{03B5}{\ensuremath{\varepsilon}}
\DeclareUnicodeCharacter{03B6}{\ensuremath{\zeta}}
\DeclareUnicodeCharacter{03B7}{\ensuremath{\eta}}
\DeclareUnicodeCharacter{03B8}{\ensuremath{\theta}}
\DeclareUnicodeCharacter{03B9}{\ensuremath{\iota}}
\DeclareUnicodeCharacter{03BA}{\ensuremath{\kappa}}
\DeclareUnicodeCharacter{03BB}{\ensuremath{\lambda}}
\DeclareUnicodeCharacter{03BC}{\ensuremath{\mu}}
\DeclareUnicodeCharacter{03BD}{\ensuremath{\nu}}
\DeclareUnicodeCharacter{03BE}{\ensuremath{\xi}}
\DeclareUnicodeCharacter{03C0}{\ensuremath{\pi}}
\DeclareUnicodeCharacter{03C1}{\ensuremath{\rho}}
\DeclareUnicodeCharacter{03C3}{\ensuremath{\sigma}}
\DeclareUnicodeCharacter{03C4}{\ensuremath{\tau}}
\DeclareUnicodeCharacter{03C5}{\ensuremath{\upsilon}}
\DeclareUnicodeCharacter{03C6}{\ensuremath{\varphi}}
\DeclareUnicodeCharacter{03C7}{\ensuremath{\chi}}
\DeclareUnicodeCharacter{03C8}{\ensuremath{\psi}}
\DeclareUnicodeCharacter{03C9}{\ensuremath{\omega}}
\DeclareUnicodeCharacter{0394}{\ensuremath{\Delta}}
\DeclareUnicodeCharacter{0398}{\ensuremath{\Theta}}
\DeclareUnicodeCharacter{039B}{\ensuremath{\Lambda}}
\DeclareUnicodeCharacter{03A0}{\ensuremath{\Pi}}
\DeclareUnicodeCharacter{03A3}{\ensuremath{\Sigma}}
\DeclareUnicodeCharacter{03A6}{\ensuremath{\Phi}}
\DeclareUnicodeCharacter{03A8}{\ensuremath{\Psi}}
\DeclareUnicodeCharacter{03A9}{\ensuremath{\Omega}}
\DeclareUnicodeCharacter{2200}{\ensuremath{\forall}}
\DeclareUnicodeCharacter{2203}{\ensuremath{\exists}}
\DeclareUnicodeCharacter{2207}{\ensuremath{\nabla}}
\DeclareUnicodeCharacter{2202}{\ensuremath{\partial}}
\DeclareUnicodeCharacter{2218}{\ensuremath{\circ}}
\DeclareUnicodeCharacter{2032}{\ensuremath{\prime}}
\DeclareUnicodeCharacter{2113}{\ensuremath{\ell}}
\DeclareUnicodeCharacter{210F}{\ensuremath{\hbar}}
\DeclareUnicodeCharacter{27E8}{\ensuremath{\langle}}
\DeclareUnicodeCharacter{27E9}{\ensuremath{\rangle}}
\DeclareUnicodeCharacter{2308}{\ensuremath{\lceil}}
\DeclareUnicodeCharacter{2309}{\ensuremath{\rceil}}
\DeclareUnicodeCharacter{230A}{\ensuremath{\lfloor}}
\DeclareUnicodeCharacter{230B}{\ensuremath{\rfloor}}
\DeclareUnicodeCharacter{222B}{\ensuremath{\int}}
\DeclareUnicodeCharacter{2211}{\ensuremath{\sum}}
\DeclareUnicodeCharacter{220F}{\ensuremath{\prod}}
\DeclareUnicodeCharacter{221E}{\ensuremath{\infty}}
\DeclareUnicodeCharacter{00B5}{\ensuremath{\mu}}
\DeclareUnicodeCharacter{2236}{\ensuremath{\colon}}
\DeclareUnicodeCharacter{2192}{\ensuremath{\to}}
\DeclareUnicodeCharacter{21A6}{\ensuremath{\mapsto}}
\DeclareUnicodeCharacter{21D2}{\ensuremath{\Rightarrow}}
\DeclareUnicodeCharacter{21D4}{\ensuremath{\Leftrightarrow}}
\DeclareUnicodeCharacter{2282}{\ensuremath{\subset}}
\DeclareUnicodeCharacter{2286}{\ensuremath{\subseteq}}
\DeclareUnicodeCharacter{2287}{\ensuremath{\supseteq}}
\DeclareUnicodeCharacter{222B}{\ensuremath{\int}}
\DeclareUnicodeCharacter{2205}{\ensuremath{\emptyset}}
\DeclareUnicodeCharacter{2299}{\ensuremath{\odot}}
\DeclareUnicodeCharacter{2220}{\ensuremath{\angle}}
\DeclareUnicodeCharacter{2135}{\ensuremath{\aleph}}
\DeclareUnicodeCharacter{207B}{\ensuremath{^{-}}}
\DeclareUnicodeCharacter{207A}{\ensuremath{^{+}}}
\DeclareUnicodeCharacter{220E}{\ensuremath{\blacksquare}}
\DeclareUnicodeCharacter{2061}{}
\DeclareUnicodeCharacter{2062}{}
\DeclareUnicodeCharacter{2063}{}
\DeclareUnicodeCharacter{2009}{\,}
\DeclareUnicodeCharacter{200A}{\,}
\DeclareUnicodeCharacter{200B}{}
\DeclareUnicodeCharacter{2212}{\ensuremath{-}}
\DeclareUnicodeCharacter{2194}{\ensuremath{\leftrightarrow}}
\DeclareUnicodeCharacter{2099}{\ensuremath{_{n}}}
\DeclareUnicodeCharacter{2098}{\ensuremath{_{m}}}
\DeclareUnicodeCharacter{2096}{\ensuremath{_{k}}}
\DeclareUnicodeCharacter{2095}{\ensuremath{_{h}}}
\DeclareUnicodeCharacter{2090}{\ensuremath{_{a}}}
\DeclareUnicodeCharacter{2091}{\ensuremath{_{e}}}
\DeclareUnicodeCharacter{2093}{\ensuremath{_{x}}}
\DeclareUnicodeCharacter{1D62}{\ensuremath{_{i}}}
\DeclareUnicodeCharacter{2C7C}{\ensuremath{_{j}}}
\DeclareUnicodeCharacter{2C7B}{\ensuremath{^{e}}}
\DeclareUnicodeCharacter{1D52}{\ensuremath{^{o}}}
\DeclareUnicodeCharacter{1D5}{\ensuremath{^{u}}}
\DeclareUnicodeCharacter{02E1}{\ensuremath{^{l}}}
\DeclareUnicodeCharacter{2295}{\ensuremath{\oplus}}
\DeclareUnicodeCharacter{27FA}{\ensuremath{\Longleftrightarrow}}
\DeclareUnicodeCharacter{27F9}{\ensuremath{\Longrightarrow}}
\DeclareUnicodeCharacter{2299}{\ensuremath{\odot}}
% Provide \tightlist used by pandoc-generated lists (no-op)
\providecommand{\tightlist}{\setlength{\itemsep}{0pt}\setlength{\parskip}{0pt}}
% pandoc longtable column specs: provide \real as passthrough and a
% `none` counter so that calc's \real{x} expansion path works.
\newcounter{none}
\providecommand{\real}[1]{#1}

% LISTINGS
\lstset{
    basicstyle=\small\ttfamily,
    breaklines=true, frame=single,
    keepspaces=true, showstringspaces=false,
    aboveskip=0.8em, belowskip=0.8em
}

% HEADER/FOOTER
\setlength{\headheight}{14pt}
\pagestyle{fancy}
\fancyhf{}
\fancyhead[R]{\thepage}
\renewcommand{\headrulewidth}{0.2pt}

% HYPERREF
\hypersetup{
    colorlinks=true,
    linkcolor=blue, citecolor=blue, urlcolor=blue,
    pdfauthor={Brieuc de La Fourniere}
}

% SPACING
\setstretch{1.2}
\setlength{\parskip}{0.4em}
\setlength{\parindent}{0pt}

% TITLE FORMATTING
\usepackage{titling}
\pretitle{\LARGE\bfseries}
\posttitle{\vspace{-0.4em}}
\preauthor{}
\postauthor{}
\predate{}
\postdate{}
\setlength{\droptitle}{-2.0em}

% CUSTOM COMMANDS — GIFT notation
\newcommand{\E}{\mathrm{E}}
\newcommand{\Gtwo}{\mathrm{G}_2}
\newcommand{\Kseven}{K_7}
\newcommand{\AdS}{\mathrm{AdS}}
\newcommand{\SU}{\mathrm{SU}}
\newcommand{\SO}{\mathrm{SO}}
\newcommand{\U}{\mathrm{U}}
\newcommand{\Sp}{\mathrm{Sp}}
\newcommand{\PSL}{\mathrm{PSL}}
\newcommand{\SM}{\mathrm{SM}}
\newcommand{\Tr}{\mathrm{Tr}}
\newcommand{\rank}{\mathrm{rank}}
\newcommand{\Vol}{\mathrm{Vol}}
\newcommand{\Hol}{\mathrm{Hol}}
\newcommand{\Ric}{\mathrm{Ric}}
\newcommand{\Rm}{\mathrm{Rm}}
\newcommand{\Scal}{\mathrm{Scal}}
\DeclareMathOperator{\diag}{diag}
\DeclareMathOperator{\cond}{cond}
\newcommand{\GIFT}{\textrm{GIFT}}
\newcommand{\CP}{\mathrm{CP}}
\newcommand{\CKM}{\mathrm{CKM}}
\newcommand{\PMNS}{\mathrm{PMNS}}
\newcommand{\EW}{\mathrm{EW}}
\newcommand{\Pl}{\mathrm{Pl}}
\newcommand{\DE}{\mathrm{DE}}
\newcommand{\DM}{\mathrm{DM}}
\newcommand{\KK}{\mathrm{KK}}
\newcommand{\NK}{\mathrm{NK}}
\newcommand{\TCS}{\mathrm{TCS}}
\newcommand{\PINN}{\mathrm{PINN}}
\newcommand{\btwo}{b_2}
\newcommand{\bthree}{b_3}
\newcommand{\typeI}{\textsc{I}}
\newcommand{\typeII}{\textsc{II}}
\newcommand{\typeIII}{\textsc{III}}
\newcommand{\typeIV}{\textsc{IV}}
\newcommand{\proven}{\textsc{Proven}}

% PDF bookmark safety
\pdfstringdefDisableCommands{%
  \def\textsubscript#1{#1}%
  \def\textsuperscript#1{#1}%
  \def\CP{CP}%
  \def\Gtwo{G2}%
  \def\Kseven{K7}%
  \def\GIFT{GIFT}%
  \def\E{E}%
  \def\SM{SM}%
}

% THEOREM ENVIRONMENTS
\newtheorem{theorem}{Theorem}[section]
\newtheorem{proposition}[theorem]{Proposition}
\newtheorem{lemma}[theorem]{Lemma}
\newtheorem{corollary}[theorem]{Corollary}
\theoremstyle{definition}
\newtheorem{definition}[theorem]{Definition}
\theoremstyle{remark}
\newtheorem{remark}[theorem]{Remark}
 after \documentclass

% ENCODING & FONTS
\usepackage[utf8]{inputenc}
\usepackage[T1]{fontenc}
\usepackage{lmodern}

% PAGE LAYOUT (Golden Ratio)
\usepackage[margin=1.618cm, top=2.618cm, bottom=2.618cm]{geometry}

% ESSENTIAL PACKAGES
\usepackage{float}
\usepackage{caption}
\usepackage{subcaption}
\usepackage{setspace}
\usepackage{fancyhdr}
\usepackage{xcolor}
\usepackage{hyperref}
\usepackage{csquotes}
\usepackage{amsmath}
\usepackage{amssymb}
\usepackage{amsthm}
\usepackage{mathtools}
\usepackage{booktabs}
\usepackage{longtable}
\usepackage{array}
\usepackage{multirow}
\usepackage{tikz}
\usepackage{graphicx}
\usepackage{listings}
\usepackage{enumitem}
% Unicode fallbacks (pandoc artifacts that survive markdown→LaTeX)
\DeclareUnicodeCharacter{00B0}{\ensuremath{^\circ}}
\DeclareUnicodeCharacter{221A}{\ensuremath{\sqrt{\,}}}
\DeclareUnicodeCharacter{2713}{\ensuremath{\checkmark}}
\DeclareUnicodeCharacter{2717}{\ensuremath{\times}}
\DeclareUnicodeCharacter{2070}{\ensuremath{^{0}}}
\DeclareUnicodeCharacter{00B9}{\ensuremath{^{1}}}
\DeclareUnicodeCharacter{00B2}{\ensuremath{^{2}}}
\DeclareUnicodeCharacter{00B3}{\ensuremath{^{3}}}
\DeclareUnicodeCharacter{2074}{\ensuremath{^{4}}}
\DeclareUnicodeCharacter{2075}{\ensuremath{^{5}}}
\DeclareUnicodeCharacter{2076}{\ensuremath{^{6}}}
\DeclareUnicodeCharacter{2077}{\ensuremath{^{7}}}
\DeclareUnicodeCharacter{2078}{\ensuremath{^{8}}}
\DeclareUnicodeCharacter{2079}{\ensuremath{^{9}}}
\DeclareUnicodeCharacter{2071}{\ensuremath{^{i}}}
\DeclareUnicodeCharacter{2080}{\ensuremath{_{0}}}
\DeclareUnicodeCharacter{2081}{\ensuremath{_{1}}}
\DeclareUnicodeCharacter{2082}{\ensuremath{_{2}}}
\DeclareUnicodeCharacter{2083}{\ensuremath{_{3}}}
\DeclareUnicodeCharacter{2084}{\ensuremath{_{4}}}
\DeclareUnicodeCharacter{2085}{\ensuremath{_{5}}}
\DeclareUnicodeCharacter{2086}{\ensuremath{_{6}}}
\DeclareUnicodeCharacter{2087}{\ensuremath{_{7}}}
\DeclareUnicodeCharacter{2088}{\ensuremath{_{8}}}
\DeclareUnicodeCharacter{2089}{\ensuremath{_{9}}}
\DeclareUnicodeCharacter{208A}{\ensuremath{_{+}}}
\DeclareUnicodeCharacter{208B}{\ensuremath{_{-}}}
\DeclareUnicodeCharacter{2192}{\ensuremath{\to}}
\DeclareUnicodeCharacter{2502}{|}
\DeclareUnicodeCharacter{25BC}{\ensuremath{\blacktriangledown}}
\DeclareUnicodeCharacter{2500}{\ensuremath{-}}
\DeclareUnicodeCharacter{2514}{\ensuremath{\llcorner}}
\DeclareUnicodeCharacter{251C}{\ensuremath{\vdash}}
\DeclareUnicodeCharacter{2026}{\ldots}
\DeclareUnicodeCharacter{2208}{\ensuremath{\in}}
\DeclareUnicodeCharacter{2229}{\ensuremath{\cap}}
\DeclareUnicodeCharacter{222A}{\ensuremath{\cup}}
\DeclareUnicodeCharacter{2295}{\ensuremath{\oplus}}
\DeclareUnicodeCharacter{2297}{\ensuremath{\otimes}}
\DeclareUnicodeCharacter{2245}{\ensuremath{\cong}}
\DeclareUnicodeCharacter{2248}{\ensuremath{\approx}}
\DeclareUnicodeCharacter{2260}{\ensuremath{\neq}}
\DeclareUnicodeCharacter{2261}{\ensuremath{\equiv}}
\DeclareUnicodeCharacter{2264}{\ensuremath{\leq}}
\DeclareUnicodeCharacter{2265}{\ensuremath{\geq}}
\DeclareUnicodeCharacter{00D7}{\ensuremath{\times}}
\DeclareUnicodeCharacter{00B1}{\ensuremath{\pm}}
\DeclareUnicodeCharacter{2225}{\ensuremath{\|}}
% Greek letters used in body text (pandoc passes them through unchanged)
\DeclareUnicodeCharacter{03B1}{\ensuremath{\alpha}}
\DeclareUnicodeCharacter{03B2}{\ensuremath{\beta}}
\DeclareUnicodeCharacter{03B3}{\ensuremath{\gamma}}
\DeclareUnicodeCharacter{03B4}{\ensuremath{\delta}}
\DeclareUnicodeCharacter{03B5}{\ensuremath{\varepsilon}}
\DeclareUnicodeCharacter{03B6}{\ensuremath{\zeta}}
\DeclareUnicodeCharacter{03B7}{\ensuremath{\eta}}
\DeclareUnicodeCharacter{03B8}{\ensuremath{\theta}}
\DeclareUnicodeCharacter{03B9}{\ensuremath{\iota}}
\DeclareUnicodeCharacter{03BA}{\ensuremath{\kappa}}
\DeclareUnicodeCharacter{03BB}{\ensuremath{\lambda}}
\DeclareUnicodeCharacter{03BC}{\ensuremath{\mu}}
\DeclareUnicodeCharacter{03BD}{\ensuremath{\nu}}
\DeclareUnicodeCharacter{03BE}{\ensuremath{\xi}}
\DeclareUnicodeCharacter{03C0}{\ensuremath{\pi}}
\DeclareUnicodeCharacter{03C1}{\ensuremath{\rho}}
\DeclareUnicodeCharacter{03C3}{\ensuremath{\sigma}}
\DeclareUnicodeCharacter{03C4}{\ensuremath{\tau}}
\DeclareUnicodeCharacter{03C5}{\ensuremath{\upsilon}}
\DeclareUnicodeCharacter{03C6}{\ensuremath{\varphi}}
\DeclareUnicodeCharacter{03C7}{\ensuremath{\chi}}
\DeclareUnicodeCharacter{03C8}{\ensuremath{\psi}}
\DeclareUnicodeCharacter{03C9}{\ensuremath{\omega}}
\DeclareUnicodeCharacter{0394}{\ensuremath{\Delta}}
\DeclareUnicodeCharacter{0398}{\ensuremath{\Theta}}
\DeclareUnicodeCharacter{039B}{\ensuremath{\Lambda}}
\DeclareUnicodeCharacter{03A0}{\ensuremath{\Pi}}
\DeclareUnicodeCharacter{03A3}{\ensuremath{\Sigma}}
\DeclareUnicodeCharacter{03A6}{\ensuremath{\Phi}}
\DeclareUnicodeCharacter{03A8}{\ensuremath{\Psi}}
\DeclareUnicodeCharacter{03A9}{\ensuremath{\Omega}}
\DeclareUnicodeCharacter{2200}{\ensuremath{\forall}}
\DeclareUnicodeCharacter{2203}{\ensuremath{\exists}}
\DeclareUnicodeCharacter{2207}{\ensuremath{\nabla}}
\DeclareUnicodeCharacter{2202}{\ensuremath{\partial}}
\DeclareUnicodeCharacter{2218}{\ensuremath{\circ}}
\DeclareUnicodeCharacter{2032}{\ensuremath{\prime}}
\DeclareUnicodeCharacter{2113}{\ensuremath{\ell}}
\DeclareUnicodeCharacter{210F}{\ensuremath{\hbar}}
\DeclareUnicodeCharacter{27E8}{\ensuremath{\langle}}
\DeclareUnicodeCharacter{27E9}{\ensuremath{\rangle}}
\DeclareUnicodeCharacter{2308}{\ensuremath{\lceil}}
\DeclareUnicodeCharacter{2309}{\ensuremath{\rceil}}
\DeclareUnicodeCharacter{230A}{\ensuremath{\lfloor}}
\DeclareUnicodeCharacter{230B}{\ensuremath{\rfloor}}
\DeclareUnicodeCharacter{222B}{\ensuremath{\int}}
\DeclareUnicodeCharacter{2211}{\ensuremath{\sum}}
\DeclareUnicodeCharacter{220F}{\ensuremath{\prod}}
\DeclareUnicodeCharacter{221E}{\ensuremath{\infty}}
\DeclareUnicodeCharacter{00B5}{\ensuremath{\mu}}
\DeclareUnicodeCharacter{2236}{\ensuremath{\colon}}
\DeclareUnicodeCharacter{2192}{\ensuremath{\to}}
\DeclareUnicodeCharacter{21A6}{\ensuremath{\mapsto}}
\DeclareUnicodeCharacter{21D2}{\ensuremath{\Rightarrow}}
\DeclareUnicodeCharacter{21D4}{\ensuremath{\Leftrightarrow}}
\DeclareUnicodeCharacter{2282}{\ensuremath{\subset}}
\DeclareUnicodeCharacter{2286}{\ensuremath{\subseteq}}
\DeclareUnicodeCharacter{2287}{\ensuremath{\supseteq}}
\DeclareUnicodeCharacter{222B}{\ensuremath{\int}}
\DeclareUnicodeCharacter{2205}{\ensuremath{\emptyset}}
\DeclareUnicodeCharacter{2299}{\ensuremath{\odot}}
\DeclareUnicodeCharacter{2220}{\ensuremath{\angle}}
\DeclareUnicodeCharacter{2135}{\ensuremath{\aleph}}
\DeclareUnicodeCharacter{207B}{\ensuremath{^{-}}}
\DeclareUnicodeCharacter{207A}{\ensuremath{^{+}}}
\DeclareUnicodeCharacter{220E}{\ensuremath{\blacksquare}}
\DeclareUnicodeCharacter{2061}{}
\DeclareUnicodeCharacter{2062}{}
\DeclareUnicodeCharacter{2063}{}
\DeclareUnicodeCharacter{2009}{\,}
\DeclareUnicodeCharacter{200A}{\,}
\DeclareUnicodeCharacter{200B}{}
\DeclareUnicodeCharacter{2212}{\ensuremath{-}}
\DeclareUnicodeCharacter{2194}{\ensuremath{\leftrightarrow}}
\DeclareUnicodeCharacter{2099}{\ensuremath{_{n}}}
\DeclareUnicodeCharacter{2098}{\ensuremath{_{m}}}
\DeclareUnicodeCharacter{2096}{\ensuremath{_{k}}}
\DeclareUnicodeCharacter{2095}{\ensuremath{_{h}}}
\DeclareUnicodeCharacter{2090}{\ensuremath{_{a}}}
\DeclareUnicodeCharacter{2091}{\ensuremath{_{e}}}
\DeclareUnicodeCharacter{2093}{\ensuremath{_{x}}}
\DeclareUnicodeCharacter{1D62}{\ensuremath{_{i}}}
\DeclareUnicodeCharacter{2C7C}{\ensuremath{_{j}}}
\DeclareUnicodeCharacter{2C7B}{\ensuremath{^{e}}}
\DeclareUnicodeCharacter{1D52}{\ensuremath{^{o}}}
\DeclareUnicodeCharacter{1D5}{\ensuremath{^{u}}}
\DeclareUnicodeCharacter{02E1}{\ensuremath{^{l}}}
\DeclareUnicodeCharacter{2295}{\ensuremath{\oplus}}
\DeclareUnicodeCharacter{27FA}{\ensuremath{\Longleftrightarrow}}
\DeclareUnicodeCharacter{27F9}{\ensuremath{\Longrightarrow}}
\DeclareUnicodeCharacter{2299}{\ensuremath{\odot}}
% Provide \tightlist used by pandoc-generated lists (no-op)
\providecommand{\tightlist}{\setlength{\itemsep}{0pt}\setlength{\parskip}{0pt}}
% pandoc longtable column specs: provide \real as passthrough and a
% `none` counter so that calc's \real{x} expansion path works.
\newcounter{none}
\providecommand{\real}[1]{#1}

% LISTINGS
\lstset{
    basicstyle=\small\ttfamily,
    breaklines=true, frame=single,
    keepspaces=true, showstringspaces=false,
    aboveskip=0.8em, belowskip=0.8em
}

% HEADER/FOOTER
\setlength{\headheight}{14pt}
\pagestyle{fancy}
\fancyhf{}
\fancyhead[R]{\thepage}
\renewcommand{\headrulewidth}{0.2pt}

% HYPERREF
\hypersetup{
    colorlinks=true,
    linkcolor=blue, citecolor=blue, urlcolor=blue,
    pdfauthor={Brieuc de La Fourniere}
}

% SPACING
\setstretch{1.2}
\setlength{\parskip}{0.4em}
\setlength{\parindent}{0pt}

% TITLE FORMATTING
\usepackage{titling}
\pretitle{\LARGE\bfseries}
\posttitle{\vspace{-0.4em}}
\preauthor{}
\postauthor{}
\predate{}
\postdate{}
\setlength{\droptitle}{-2.0em}

% CUSTOM COMMANDS — GIFT notation
\newcommand{\E}{\mathrm{E}}
\newcommand{\Gtwo}{\mathrm{G}_2}
\newcommand{\Kseven}{K_7}
\newcommand{\AdS}{\mathrm{AdS}}
\newcommand{\SU}{\mathrm{SU}}
\newcommand{\SO}{\mathrm{SO}}
\newcommand{\U}{\mathrm{U}}
\newcommand{\Sp}{\mathrm{Sp}}
\newcommand{\PSL}{\mathrm{PSL}}
\newcommand{\SM}{\mathrm{SM}}
\newcommand{\Tr}{\mathrm{Tr}}
\newcommand{\rank}{\mathrm{rank}}
\newcommand{\Vol}{\mathrm{Vol}}
\newcommand{\Hol}{\mathrm{Hol}}
\newcommand{\Ric}{\mathrm{Ric}}
\newcommand{\Rm}{\mathrm{Rm}}
\newcommand{\Scal}{\mathrm{Scal}}
\DeclareMathOperator{\diag}{diag}
\DeclareMathOperator{\cond}{cond}
\newcommand{\GIFT}{\textrm{GIFT}}
\newcommand{\CP}{\mathrm{CP}}
\newcommand{\CKM}{\mathrm{CKM}}
\newcommand{\PMNS}{\mathrm{PMNS}}
\newcommand{\EW}{\mathrm{EW}}
\newcommand{\Pl}{\mathrm{Pl}}
\newcommand{\DE}{\mathrm{DE}}
\newcommand{\DM}{\mathrm{DM}}
\newcommand{\KK}{\mathrm{KK}}
\newcommand{\NK}{\mathrm{NK}}
\newcommand{\TCS}{\mathrm{TCS}}
\newcommand{\PINN}{\mathrm{PINN}}
\newcommand{\btwo}{b_2}
\newcommand{\bthree}{b_3}
\newcommand{\typeI}{\textsc{I}}
\newcommand{\typeII}{\textsc{II}}
\newcommand{\typeIII}{\textsc{III}}
\newcommand{\typeIV}{\textsc{IV}}
\newcommand{\proven}{\textsc{Proven}}

% PDF bookmark safety
\pdfstringdefDisableCommands{%
  \def\textsubscript#1{#1}%
  \def\textsuperscript#1{#1}%
  \def\CP{CP}%
  \def\Gtwo{G2}%
  \def\Kseven{K7}%
  \def\GIFT{GIFT}%
  \def\E{E}%
  \def\SM{SM}%
}

% THEOREM ENVIRONMENTS
\newtheorem{theorem}{Theorem}[section]
\newtheorem{proposition}[theorem]{Proposition}
\newtheorem{lemma}[theorem]{Lemma}
\newtheorem{corollary}[theorem]{Corollary}
\theoremstyle{definition}
\newtheorem{definition}[theorem]{Definition}
\theoremstyle{remark}
\newtheorem{remark}[theorem]{Remark}

\fancyhead[L]{GIFT v3.4: Supplement S3}
\hypersetup{pdftitle={GIFT v3.4 Supplement S3: Observable Dataset}}

\title{%
\LARGE\textbf{Supplement S3: Observable Dataset}\\[0.3em]
\large Complete GIFT Observable Dataset (95 entries)%
}
\author{Brieuc de La Fourni\`ere \\ \textit{Independent researcher, Beaune, France}}
\date{May 2026 \quad\textbar\quad v3.4}

\graphicspath{{../../../figures/}}

\begin{document}

\maketitle
\noindent\rule{\textwidth}{0.2pt}
\thispagestyle{fancy}

\tableofcontents
\newpage


\subsection{S3.1 Introduction}\label{s3.1-introduction}

This supplement provides the complete GIFT observable dataset: 95 observables derived from a single G₂ metric (169 Chebyshev-Cholesky geometric parameters, zero freely adjustable physical parameters), organized by derivation type and cross-referenced with Lean formal verification. The dataset includes 2 BH remnant predictions from Pinčák et al.~2026 \cite{46} (classified Type IV) and 6 new Type IV metric block eigenvalues.

The 95 observables decompose as \textbf{33(I) + 19(II) + 21(III) + 22(IV)} across four types of increasing derivation complexity. All predictions trace to a single compact G₂ manifold K₇ with (b₂, b₃) = (21, 77) and E₈×E₈ gauge structure, 169 metric parameters, zero continuously adjustable physical parameters. The pair (b₂, b₃) = (21, 77) does not appear among the catalogued compact G₂ manifolds; a complete geometric construction remains an open problem. The certified metric provides computational evidence for a new G₂ manifold.

The dataset serves as the canonical reference for all GIFT predictions, superseding scattered results across supplements S6--S25. Every entry traces back to a specific computation, a JSON artifact, and (where applicable) a Lean theorem.

\begin{center}\rule{0.5\linewidth}{0.5pt}\end{center}

\subsection{S3.2 Type Classification}\label{s3.2-type-classification}

Observables are classified into four types by derivation directness. The type assignment reflects the number of physical identification steps between the G₂ metric and the final prediction:

{\def\LTcaptype{none} % do not increment counter
\begin{longtable}[]{@{}
  >{\raggedright\arraybackslash}p{(\linewidth - 8\tabcolsep) * \real{0.1017}}
  >{\raggedright\arraybackslash}p{(\linewidth - 8\tabcolsep) * \real{0.1186}}
  >{\raggedright\arraybackslash}p{(\linewidth - 8\tabcolsep) * \real{0.2203}}
  >{\raggedright\arraybackslash}p{(\linewidth - 8\tabcolsep) * \real{0.3051}}
  >{\raggedright\arraybackslash}p{(\linewidth - 8\tabcolsep) * \real{0.2542}}@{}}
\toprule\noalign{}
\begin{minipage}[b]{\linewidth}\raggedright
Type
\end{minipage} & \begin{minipage}[b]{\linewidth}\raggedright
Count
\end{minipage} & \begin{minipage}[b]{\linewidth}\raggedright
Description
\end{minipage} & \begin{minipage}[b]{\linewidth}\raggedright
Derivation chain
\end{minipage} & \begin{minipage}[b]{\linewidth}\raggedright
Lean coverage
\end{minipage} \\
\midrule\noalign{}
\endhead
\bottomrule\noalign{}
\endlastfoot
I & 33 & Direct algebraic from structural constants & topology → formula → ratio & 33/33 (100\%) \\
II & 19 & One physical identification step & formula + VEV or scale → dimensional quantity & 0/19 (0\%) \\
III & 21 & Multi-step dynamical chains & metric → mechanism → observable & 14/21 (67\%) \\
IV & 22 & Structural/internal quantities & metric → diagnostic & 8/22 (36\%) \\
\end{longtable}
\textbf{Type I} (33 observables): Direct from metric geometry and topological invariants (b₂, b₃, dim(G₂), dim(E₈), etc.). These are dimensionless ratios expressed as algebraic combinations of integers. Examples: sin²θ\_W = 3/13, Q\_Koide = 2/3, m\_τ/m\_e = 3477. All 33 are Lean-certified. The two BH remnant topological predictions are classified as Type IV structural diagnostics (not part of the 33 Type I Lean-certified core).

\textbf{Type II} (19 observables): One physical identification step beyond Type I. Absolute masses from ratios × VEV (e.g., m\_u = (m\_u/m\_d) × m\_d\_exp), CKM magnitudes from Wolfenstein parametrization, extended quark ratios. From the \texttt{gift\_observables.csv} pipeline.

\textbf{Type III} (21 observables): Multi-step chains involving dynamical mechanisms:
- wilson\_line non-adiabatic eigenvalue splitting (3 observables: raw lepton ratios)
- S11 RGE gauge coupling running (4 observables: sin²θ\_W, α\_em⁻¹, α\_s at M\_Z, M\_GUT)
- spectral 7D spectral geometry (5 observables: Weyl exponent, KK states, fiber channels)
- gauge\_bundle gauge bundle data (4 observables: f\_IJ conditioning, α\_ratio, Yukawa rank)
- instanton instanton volumes + combined wilson\_line+instanton (5 observables: ΔV's, combined ratios)

\textbf{Type IV} (22 observables): Structural and internal consistency quantities with no direct experimental comparison. These include topology diagnostics (b₂, b₃, χ(K₇)), Newton-Kantorovich certification values (h, distance, margin), Gram matrix conditioning (cond(G\_K3), cond(G\_77)), spectral counts, metric block eigenvalues (g\_ss = 19/6, g\_\{T²\} = 7/6, g\_K3 ≈ 64/77), and Pinčák et al.~2026 \cite{46} topological diagnostics (N\_QNM = 98, M\_res, instanton suppression ×77). M\_res and N\_QNM are the two BH-remnant-related entries; they are classified here as structural diagnostics, not as Type I or Type III, because they require a physical scale identification (τ₀ = v\_EW) not derivable from topology alone. These quantities verify internal consistency rather than predict measurements directly.

\textbf{Derivation depth}: The type assignment correlates with derivation chain length. Type I observables require 1 algebraic step; Type II requires 2 steps (algebra + scale identification); Type III requires 3--5 steps (algebra + dynamics + calibration); Type IV involves diagnostic computation chains of variable length.

\begin{center}\rule{0.5\linewidth}{0.5pt}\end{center}

\subsection{S3.3 Methodology}\label{s3.3-methodology}

\subsubsection{Experimental Sources}\label{experimental-sources}

All experimental reference values are drawn from four standard compilations:

{\def\LTcaptype{none} % do not increment counter
\begin{longtable}[]{@{}
  >{\raggedright\arraybackslash}p{(\linewidth - 4\tabcolsep) * \real{0.3333}}
  >{\raggedright\arraybackslash}p{(\linewidth - 4\tabcolsep) * \real{0.4167}}
  >{\raggedright\arraybackslash}p{(\linewidth - 4\tabcolsep) * \real{0.2500}}@{}}
\toprule\noalign{}
\begin{minipage}[b]{\linewidth}\raggedright
Source
\end{minipage} & \begin{minipage}[b]{\linewidth}\raggedright
Coverage
\end{minipage} & \begin{minipage}[b]{\linewidth}\raggedright
Date
\end{minipage} \\
\midrule\noalign{}
\endhead
\bottomrule\noalign{}
\endlastfoot
\textbf{PDG 2024} \cite{1} & Particle masses, gauge couplings, CKM elements & 2025 \\
\textbf{NuFIT 6.0} \cite{2} & Neutrino oscillation parameters (NO w/o SK) & Oct 2024 \\
\textbf{Planck PR4} (Tristram+ 2024) \cite{3} & Cosmological parameters: h, σ₈, Ω\_Λ & 2024 \\
\textbf{CODATA 2022} \cite{4} & Lepton mass ratios & 2022 \\
\end{longtable}
Selected NuFIT 6.0 values: sin²θ₂₃ = 0.561, δ\_CP = 177°. Selected Planck PR4 values: h = 0.6764, σ₈ = 0.807, Ω\_Λ = 0.6847.

\subsubsection{Structural Constants}\label{structural-constants}

The 20 structural constants from which all 95 observables derive:

{\def\LTcaptype{none} % do not increment counter
\begin{longtable}[]{@{}
  >{\raggedright\arraybackslash}p{(\linewidth - 6\tabcolsep) * \real{0.1071}}
  >{\raggedright\arraybackslash}p{(\linewidth - 6\tabcolsep) * \real{0.3571}}
  >{\raggedright\arraybackslash}p{(\linewidth - 6\tabcolsep) * \real{0.2500}}
  >{\raggedright\arraybackslash}p{(\linewidth - 6\tabcolsep) * \real{0.2857}}@{}}
\toprule\noalign{}
\begin{minipage}[b]{\linewidth}\raggedright
\#
\end{minipage} & \begin{minipage}[b]{\linewidth}\raggedright
Constant
\end{minipage} & \begin{minipage}[b]{\linewidth}\raggedright
Value
\end{minipage} & \begin{minipage}[b]{\linewidth}\raggedright
Origin
\end{minipage} \\
\midrule\noalign{}
\endhead
\bottomrule\noalign{}
\endlastfoot
1 & dim(E₈) & 248 & E₈ Lie algebra \\
2 & rank(E₈) & 8 & Cartan subalgebra \\
3 & dim(G₂) & 14 & Aut($\mathbb{O}$) \\
4 & dim(K₇) & 7 & Im($\mathbb{O}$) \\
5 & b₂(K₇) & 21 & Input; confirmed by spectral analysis (Paper B [B], §2.3). Any Mayer-Vietoris decomposition is conditional on building block identification (open problem). \\
6 & b₃(K₇) & 77 & Input; confirmed by spectral analysis (Paper B [B], §2.3). Any Mayer-Vietoris decomposition is conditional on building block identification (open problem). \\
7 & N\_gen & 3 & Index theorem \\
8 & p₂ & 2 & dim(G₂)/dim(K₇) \\
9 & Weyl & 5 & Triple identity \\
10 & H* & 99 & b₂+b₃+1 \\
11 & dim(J₃($\mathbb{O}$)) & 27 & Jordan algebra \\
12 & dim(E₆) & 78 & E₆ Lie algebra \\
13 & dim(E₇) & 133 & E₇ Lie algebra \\
14 & dim(F₄) & 52 & Aut(J₃($\mathbb{O}$)) \\
15 & fund(E₇) & 56 & E₇ fundamental \\
16 & \textbar PSL(2,7)\textbar{} & 168 & Fano automorphisms \\
17 & D\_bulk & 11 & 4+7 \\
18 & α\_sum & 13 & rank(E₈)+Weyl \\
19 & det(g)\_den & 32 & b₂+dim(G₂)−N\_gen \\
20 & det(g)\_num & 65 & Weyl×α\_sum \\
\end{longtable}
\subsubsection{Consolidation Pipeline}\label{consolidation-pipeline}

The master dataset (\texttt{gift\_observables.json}) is generated by \texttt{run\_observable\_dataset.py} (observable\_dataset), which collects all individual JSON results from S6--S25 and performs 5 verification checks:

\begin{enumerate}


\item
  \textbf{Completeness}: All 95 entries present across 4 types
\item
  \textbf{Uniqueness}: No duplicate observable IDs
\item
  \textbf{Consistency}: GIFT values match source JSONs to machine precision
\item
  \textbf{Source tracing}: Every entry maps to a computation step (S6--S25)
\item
  \textbf{Format validation}: JSON, CSV, and LaTeX outputs agree
\end{enumerate}

Each observable record contains: GIFT prediction (float64), GIFT fraction (exact symbolic), experimental value and uncertainty, experimental source, relative deviation (\%), Lean theorem name (if certified), and provenance (computation step).

\begin{center}\rule{0.5\linewidth}{0.5pt}\end{center}

\subsection{S3.4 Complete Observable Table}\label{s3.4-complete-observable-table}

\textbf{Table S3.1}: Complete GIFT observable dataset (95 entries). Columns: name, GIFT prediction, experimental value, relative deviation (\%), type (I/II/III/IV), source computation (S6--S25), Lean verification status.

{[}The full LaTeX table is provided in \texttt{data/gift\_observables\_latex.tex} and will be included at typesetting.{]}

\subsubsection{Sector-by-Sector Performance}\label{sector-by-sector-performance}

The 66 observables with experimental comparison (Types I+II+III) organize into 11 sectors with varying precision:

\textbf{Best-performing sectors} (mean dev \textless{} 0.5\%):
- \textbf{Boson} (3 obs): mean 0.13\%, m\_H/m\_W (0.02\%), m\_W/m\_Z (0.06\%), m\_H/m\_t (0.31\%)
- \textbf{Quark} (4 obs): mean 0.21\%, m\_s/m\_d (0.00\%), m\_u/m\_d (0.05\%), m\_c/m\_s (0.12\%), m\_b/m\_t (0.79\%)
- \textbf{CKM} (3 obs): mean 0.59\%: A\_Wolf (0.29\%), sin²θ₁₂ (0.36\%), sin²θ₂₃ (1.13\%)
- \textbf{Cosmology} (7 obs): mean 0.15\%, n\_s (0.004\%), h (0.09\%), Ω\_b/Ω\_m (0.16\%), σ₈ (0.18\%), Ω\_DE (0.21\%), Y\_p (0.37\%), Ω\_DM/Ω\_b (0.00\%)

\textbf{Good sectors} (mean dev 0.5--3\%):
- \textbf{Electroweak} (3 obs): mean 0.11\%, α⁻¹ (0.002\%), α\_s (0.126\%), sin²θ\_W (0.20\%)
- \textbf{Lepton} (3 obs): mean 0.04\%: Q\_Koide (0.001\%), m\_τ/m\_e (0.004\%), m\_μ/m\_e (0.12\%)
- \textbf{PMNS} (4 obs): mean 0.29\%, θ₁₂ (0.03\%), θ₂₃ (0.10\%), sin²θ₁₂ (0.23\%), θ₁₃ (0.37\%)

\textbf{Challenging sectors} (mean dev \textgreater{} 3\%):
- \textbf{Gauge running} (4 obs): mean 2.3\%, sin²θ\_W(RGE) (2.78\%), α\_em⁻¹(RGE) (2.53\%), α\_s(RGE) (3.7\%), M\_GUT (0.00\%)
- \textbf{Instanton} (3 obs): mean 7.4\%: ΔV(e-τ) (5.9\%), ΔV(e-μ) (15.9\%), κ(gauge) (4.7\%)

\textbf{B-test consistency} (derived, not counted as independent observable): The MSSM B parameter B = (α₁⁻¹−α₂⁻¹)/(α₂⁻¹−α₃⁻¹) evaluates to 1.4033 using GIFT couplings (0.23\% from the theoretical 7/5), closer than the purely experimental B = 1.3948 (0.37\%). The identity B = 7/5 is equivalent to α\_em⁻¹(M\_Z) = 91√2 = dim(Λ²$\mathfrak{g}$₂)·√2 and implies the holonomy sequence α₁⁻¹:α₂⁻¹:α₃⁻¹ = dim(G₂):dim(K₇):p₂ = 14:7:2. See §5.4 of the main text.

\subsubsection{Exact Matches}\label{exact-matches}

11 observables achieve deviations below 0.01\% (effectively exact):

{\def\LTcaptype{none} % do not increment counter
\begin{longtable}[]{@{}llll@{}}
\toprule\noalign{}
Observable & GIFT & Exp. & Dev. \\
\midrule\noalign{}
\endhead
\bottomrule\noalign{}
\endlastfoot
α⁻¹ & 137.033 & 137.036 & 0.002\% \\
Q\_Koide & 2/3 & 0.666661 & 0.001\% \\
m\_τ/m\_e & 3477 & 3477.15 & 0.004\% \\
m\_s/m\_d & 20 & 20.0 & 0.000\% \\
n\_s & 0.9649 & 0.9649 & 0.004\% \\
m\_u & 2.16 & 2.16 & 0.00\% \\
\textbar V\_tb\textbar{} & 0.999 & 0.999 & 0.00\% \\
m\_c/m\_d & 234.3 & 234.0 & 0.004\% \\
M\_GUT & 2×10¹⁶ & 2×10¹⁶ & 0.00\% \\
α\_GUT⁻¹ & 25.3 & 25.3 & 0.00\% \\
rank(Y) & 3 & 3 & 0.00\% \\
\end{longtable}
Note: δ\_CP canonical prediction is 197° (11.3\% from NuFIT 6.0 central); a compactification factor 62/69 brings it to 177.014° (0.008\%) but is not adopted, see §4.4.1. m\_c/m\_s deviation is 0.12\%, above the 0.01\% threshold.

\subsubsection{Known Outliers}\label{known-outliers}

Two observables deviate by more than 5\%, each with an identified physical origin:

\textbf{δ\_CP}: The GIFT prediction is 197° (pure topological: 7 × 14 + 99), deviating 11.3\% from NuFIT 6.0 central (177° ± 20°) but at the edge of 1σ. PSLQ analysis identifies a compactification factor 62/69 (§4.4.1 of main paper), documented as a structural observation but not adopted: the canonical prediction remains 197°. The experimental central value shifted from 197° (NuFIT 5.2) to 177° (NuFIT 6.0) and carries ±20° uncertainty. Falsifiable by DUNE (2028--2040).

\textbf{α\_s(RGE) = +3.7\%}: G₂-MSSM split-spectrum matching (MSSM above M\_squark = 3165 GeV, SM + gauginos below (b₃ = −5)) gives α\_s = 0.1224 (exp 0.1180, dev 3.7\%). A naive degenerate-spectrum treatment would give 12.1\%. The topological value α\_s = √2/12 = 0.11785 (Type I) matches experiment to 0.126\%.

\textbf{ΔV(e-μ) = +15.9\%}: The instanton volume assignment ΔV(e-μ) = 3.271 versus target ln(16.82) = 2.823 reflects the optimal assignment problem for 57 associative cycles. The combined wilson\_line+instanton pipeline with α = e\^{}K (§S3.5) reduces this to 0.75\%.

\textbf{κ(gauge) = +4.7\%}: The gauge kinetic function conditioning cond(f\_IJ) = 1.047 measures departure from exact universality. The deviation from 1.000 reflects K₇ geometry at the percent level.

\subsubsection{Summary Statistics}\label{summary-statistics}

{\def\LTcaptype{none} % do not increment counter
\begin{longtable}[]{@{}ll@{}}
\toprule\noalign{}
Metric & Value \\
\midrule\noalign{}
\endhead
\bottomrule\noalign{}
\endlastfoot
Total observables & 95 \\
Type breakdown & 33(I) + 19(II) + 21(III) + 22(IV) \\
With experimental comparison & 66 \\
Exact matches (dev \textless{} 0.01\%) & 11 \\
Within 1\% of experiment & 53 \\
Within 5\% of experiment & 63 \\
Lean-certified & 55 \\
Structural constants & 20 \\
Metric parameters & 169 \\
\end{longtable}
\begin{center}\rule{0.5\linewidth}{0.5pt}\end{center}

\subsection{S3.5 The Combined wilson\_line+instanton Pipeline}\label{s3.5-the-combined-wilson_lineinstanton-pipeline}

The lepton mass hierarchy (combined lepton ratio observables) uses two complementary geometric mechanisms. The complete derivation is in §6 of the main text; results are summarised here for cross-reference.

Raw wilson\_line (K3 fiber eigenvalue splitting at c = 0.452): m\_τ/m\_e = 3403 (2.1\%), m\_τ/m\_μ = 16.54 (1.7\%), m\_μ/m\_e = 205.7 (0.5\%). Raw instanton (associative 3-cycle volumes, 57 cycles): ΔV(e-τ) = 8.633 (5.9\%), ΔV(e-μ) = 3.271 (15.9\%). Combined with α = e\^{}\{K₀\} = $\hat{V}$\^{}\{−3\} = 0.002763 (Kähler potential of K₇, zero free parameters):

\begin{itemize}

\item
  m\_τ/m\_e = 3485 (exp 3477, \textbf{0.24\%})
\item
  m\_τ/m\_μ = 16.69 (exp 16.82, \textbf{0.75\%})
\item
  m\_μ/m\_e = 208.8 (exp 206.8, \textbf{0.97\%})
\end{itemize}

Certified in Lean: \texttt{AssociativeVolumes.lean} (14 conjuncts, 0 sorry). Full mechanism description: main §6.4.

\begin{center}\rule{0.5\linewidth}{0.5pt}\end{center}

\subsection{S3.6 Sensitivity Cross-Reference}\label{s3.6-sensitivity-cross-reference}

The complete sensitivity analysis is in §7 of the main text. Key results for cross-reference:

\begin{itemize}

\item
  \textbf{Effective DOF}: r\_eff = 15.53 (SVD of 20×33 constant-usage Jacobian)
\item
  \textbf{Overdetermination}: 33/15.53 = 2.13× (more constraints than free dimensions)
\item
  \textbf{Cross-coupling}: 155/528 observable pairs share a structural constant; all form one connected component
\item
  \textbf{P(coincidence, uniform)} = 10⁻³⁴⁶ (χ² + Fisher combined, main §7.5)
\item
  \textbf{P(coincidence, algebraic)} = 10⁻¹³³ (4.2M random formulas from same 20 constants, main §7.5)
\end{itemize}

Figures: \texttt{fig\_sensitivity\_heatmap.png}, \texttt{fig\_constant\_usage.png}, \texttt{fig\_observable\_correlations.png}, \texttt{fig\_mc\_per\_observable.png}.

\begin{center}\rule{0.5\linewidth}{0.5pt}\end{center}

\subsection{S3.7 Lean Cross-References}\label{s3.7-lean-cross-references}

Of the 95 observables, \textbf{55 are formally verified} in Lean 4 across the following certificate files:

\begin{small}
\noindent\begin{tabular*}{\linewidth}{@{\extracolsep{\fill}} l c c c l @{}}
\toprule
\textbf{Lean File} & \textbf{Axioms} & \textbf{Theorems} & \textbf{Conjuncts} & \textbf{Coverage} \\
\midrule
\texttt{Foundations.lean} & 0* & n/a & 38 & Type I: metric, torsion, topology \\
\texttt{Predictions.lean} & 0* & n/a & 55 & Types I--III: couplings, masses, mixing \\
\texttt{Spectral.lean} & 0* & n/a & 41 & Type II: KK spectrum, Weyl law \\
\texttt{MetricEigenvalues.lean} & 0 & n/a & 15 & Metric fractions \\
\texttt{SpectralInvariants.lean} & 0 & n/a & 10 & Heat kernel, $\zeta'(0)$, $b_1=0$ \\
\texttt{CompactificationCorrection.lean} & 0 & n/a & 6 & $\delta_{\mathrm{CP}}$ compactification factor \\
\texttt{TCSGaugeBreaking.lean} & 0 & 14 & 10 & Type III: gauge breaking chain \\
\texttt{GaugeBundleData.lean} & 0 & 12 & 11 & Type III: bundle universality \\
\texttt{AssociativeVolumes.lean} & 0 & 19 & 14 & Type III: instanton hierarchy \\
\texttt{ComputedWeylLaw.lean} & 0 & 8 & 7 & Type III: 7D spectral geometry \\
\bottomrule
\end{tabular*}
\end{small}
*4 published axioms. All substantive: standard theorems (Cheeger inequality) + geometric structure (TCS spectral bounds) + physical inputs (literature package: CGN+Joyce). KK\_YM\_EFT, K7\_spectral\_data, K7\_analysis\_data are theorems/defs. G₂ group structure proven by \texttt{native\_decide}.

\textbf{Total certificate}: 213 conjuncts, 0 \texttt{sorry}, 134 .lean files (128 core + 6 generated + 12 test/support), 8378 build jobs (Lean 4.29.0).

\subsubsection{Coverage by Type}\label{coverage-by-type}

{\def\LTcaptype{none} % do not increment counter
\begin{longtable}[]{@{}llll@{}}
\toprule\noalign{}
Type & Certified & Total & Coverage \\
\midrule\noalign{}
\endhead
\bottomrule\noalign{}
\endlastfoot
\textbf{I} & 33 & 33 & 100\% \\
\textbf{II} & 0 & 19 & 0\% \\
\textbf{III} & 14 & 21 & 67\% \\
\textbf{IV} & 8 & 22 & 36\% \\
\textbf{Total} & \textbf{55} & \textbf{95} & \textbf{58\%} \\
\end{longtable}
\textbf{Type II uncertified}: The 19 Type II observables derive from the CSV pipeline (absolute masses from ratio × VEV, CKM magnitudes from Wolfenstein). These involve dimensional quantities and intermediate data that have not yet been axiomatized in Lean. The underlying Type I ratios from which they derive are all certified.

\textbf{Type IV partial}: 8 of 22 structural quantities are certified (Betti numbers, holonomy dimension, NK certification bounds). The remaining 14 involve spectral counts, conditioning numbers, metric eigenvalues, and BH remnant estimates that require floating-point Lean infrastructure not yet developed.

\begin{center}\rule{0.5\linewidth}{0.5pt}\end{center}

\subsection{S3.8 Data Availability}\label{s3.8-data-availability}

The complete dataset is available in three formats:

\begin{itemize}

\item
  \textbf{JSON}: \texttt{gift\_observables.json} (machine-readable, full metadata per observable)
\item
  \textbf{CSV}: \texttt{gift\_observables.csv} (tabular, 95 rows × 13 columns, for analysis)
\item
  \textbf{LaTeX}: \texttt{gift\_observables\_latex.tex} (publication-ready table)
\end{itemize}

All artifacts are archived in the Zenodo data deposit (DOI: \href{https://doi.org/10.5281/zenodo.19893371}{10.5281/zenodo.19893371}), which also contains the 9 source JSON files from S6--S25, the 10 computation scripts, and the 10 figures.

\begin{center}\rule{0.5\linewidth}{0.5pt}\end{center}

\subsection{References}\label{references}

\cite{1} Particle Data Group, Phys. Rev.~D 110, 030001 (2024)
\cite{2} NuFIT 6.0, www.nu-fit.org (Oct 2024)
\cite{3} L. Tristram et al., A\&A 682, A37 (2024) (Planck PR4)
\cite{4} CODATA 2022, physics.nist.gov/cuu/Constants

\begin{center}\rule{0.5\linewidth}{0.5pt}\end{center}

\emph{GIFT Framework: Supplement S3: Observable Dataset}
\emph{95 observables, 4 types, 55 Lean-certified}


\section*{References}
\addcontentsline{toc}{section}{References}

\begin{thebibliography}{99}
\bibitem{1} Particle Data Group, ``Review of Particle Physics,'' \textit{Phys.\ Rev.\ D} \textbf{110}, 030001 (2024).
\bibitem{2} NuFIT 6.0, \url{https://www.nu-fit.org} (Oct 2024).
\bibitem{3} L.~Tristram et al., ``Planck PR4 cosmological parameters,'' \textit{Astron.\ Astrophys.} \textbf{682}, A37 (2024).
\bibitem{4} CODATA 2022 recommended values, \url{https://physics.nist.gov/cuu/Constants}.
\bibitem{46} R.~Pin\v{c}\'ak, A.~Pigazzini, M.~Pudl\'ak, E.~Barto\v{s}, ``Geometric origin of a stable black hole remnant from torsion in $G_2$-manifold geometry,'' \textit{Gen.\ Rel.\ Grav.} \textbf{58}, 29 (2026).
\end{thebibliography}

\noindent\rule{\textwidth}{0.2pt}

\begin{center}
\emph{GIFT Framework --- Supplement S3: Observable Dataset --- 95 observables, 4 types, 55 Lean-certified}
\end{center}

\end{document}
